\section{Introduction}
\label{sec:intro}

A soft sensor that reports an uncertainty interval makes two promises, and they are
not the same promise. One is \emph{accuracy}: the point estimate is close to the true
quality variable. The other is \emph{validity}: the reported interval contains the
truth at the stated rate. A sensor can be accurate but overconfident, or inaccurate but
honestly uncertain, and an operator who cannot tell which cannot safely act on either.
Most soft-sensor practice reports a point value with an error bar derived from a
held-out set and treats the two as one quantity. They are separable, and this paper
separates them.

The separation is not merely conceptual. We build a soft sensor whose accuracy comes
from process physics and whose validity comes from distribution-free calibration, then
test the two promises independently by constructing a process on which the first is
designed to fail. On a deliberately physics-free dataset the point model collapses to
worse-than-mean prediction, yet the calibrated intervals retain their coverage. That
dissociation, accuracy gone and validity intact, is the central result, and it licenses
a claim that no positive result alone could support: that in this framework
distribution-free validity is model-agnostic, while point accuracy is purchased
specifically by the physics. A negative control is the standard instrument for
attributing an effect to its cause; it is uncommon in data-driven process modeling, and
deploying one here is the paper's primary methodological contribution.

Reaching a sensor on which that test is meaningful requires getting four ordinary things
right, because each, left unaddressed, produces a sensor that is inaccurate or invalid
for reasons unrelated to the physics, and would confound the control. Process data are
autocorrelated and regime-structured, so model selection on adjacent or shuffled samples
leaks information and becomes a lottery whose winner fails on the regime the validation
scheme hid. Calibration drifts between laboratory analyses, and those analyses arrive
late and sparsely, so the labels that could correct the drift are themselves delayed.
Prediction intervals inherited from a held-out set assume exchangeability, the very
property that drift destroys, so the stated coverage is wrong in both directions. And
when a delayed label finally arrives, naive online calibration scores it against the
current interval rather than the one issued when the sample was measured, silently
corrupting the adaptation meant to fix the previous problem. We treat each with the
simplest mechanism that works, and we keep every point model linear, both because the
negative-control logic requires holding model capacity fixed across datasets and because
linearity keeps each adaptive step $O(1)$ per label and the implemented sensor auditable
by a process engineer.

The contributions, ordered by how far they depart from existing practice, are:

\begin{itemize}
\tightlist
\item
  \textbf{C1, a negative control that isolates the source of the guarantee.} On SECOM,
  a deliberately physics-free process where no thermodynamic structure exists to encode,
  the point sensor is designed to fail and does ($R^2 = -1.84$), yet the calibration
  machinery holds $0.910$--$0.915$ coverage at $37\%$ narrower width than a static
  baseline. Accuracy collapses while validity survives, which isolates physics anchors
  as the source of accuracy and distribution-free calibration as the source of validity.
  To our knowledge this is the first use of a designed negative control to attribute the
  guarantee of a process soft sensor to its cause.
\item
  \textbf{C2, leakage-free model selection.} Blocked time-series cross-validation whose
  per-fold spread is treated as signal, plus the one-standard-error parsimony rule,
  applied identically to feature sets and regularization paths, converts the cross-regime
  selection lottery into reproducible choice. On the debutanizer benchmark this separates
  a 4-feature physics set (CV $R^2$ $+0.145 \pm 0.419$, test $+0.476$) from a 126-feature
  lagged set whose failure is visible even out-of-sample (test $+0.221$); on SECOM
  ($p \approx n$) it contains a path explosion spanning four orders of magnitude.
\item
  \textbf{C3, drift handling at the documented label delay.} An open-loop EWMA bias
  update cuts cross-fold standard error $9\times$ and lifts held-out $R^2$ from $+0.476$
  to $+0.857$ on the debutanizer, with drift alarming on the corrected residual so an
  alarm signals intervention rather than normal tracking.
\item
  \textbf{C4, calibrated uncertainty under drift and delay.} Adaptive conformal intervals
  on the bias-corrected estimate achieve regime-uniform coverage ($0.897$--$0.903$ against
  a $0.90$ target, all regimes $\times$ delays) where the static construction is invalid in
  both directions ($0.847$--$0.957$). Under delayed feedback the coverage indicator is
  scored against the interval issued when the sample was measured, not the interval the
  adapted state would issue at result arrival; we state this correctness property and
  enforce it in a real-time implementation that runs two orders of magnitude inside a
  typical analyzer cycle.
\item
  \textbf{C5, transfer with its limits stated.} The methodology transfers across process
  topology verbatim; within a process, functional scale-bias migration with a GP bias
  achieves $10.0\times$ data efficiency on both evaluated regime pairs, with intervals
  that widen appropriately with the regime gap, while the linear-source degeneracy of
  per-input affine migration is verified analytically and observed empirically to
  $\pm 0.001$--$0.006$.
\end{itemize}

The paper proceeds as outlined: Section 2 positions the three literatures and the gap;
Section 3 presents the framework, each component introduced by the failure mode it
removes; Section 4 specifies the three case studies and the shared protocol; Section 5
reports results; Section 6 discusses the production stance, the linear-scope tradeoff,
and limitations; Section 7 concludes.
